\documentclass[12pt]{article}
\usepackage[utf8]{inputenc}
\usepackage[T1]{fontenc}
\usepackage[margin=1in]{geometry}
\usepackage{setspace}
\usepackage{xcolor}
\usepackage[colorlinks=true, urlcolor=blue!60!black, linkcolor=blue!60!black]{hyperref}
\setlength{\parindent}{2em}
\setlength{\parskip}{1.25em}
\renewcommand{\baselinestretch}{1.0}
\begin{document}
\begin{center}
{\LARGE \textit{Bolling v. Sharpe}, 347 U.S. 497 (1954)\par}
\vspace{0.8em}
{\large Case Brief --- Four-Page Edition\par}
\end{center}
\vspace{1em}

\section{Memory Jogger}

Because it would be ``unthinkable'' for the Constitution to forbid racial school segregation by the states while allowing it by the federal government, the Fifth Amendment's Due Process Clause prohibits segregated public schools in Washington, D.C., launching the doctrine of reverse incorporation (\href{https://supreme.justia.com/cases/federal/us/347/497/}{347 U.S. 497}).

\section{Detailed Case Facts}

In 1947, D.C. barber and activist Gardner Bishop co-founded the Consolidated Parents Group to challenge overcrowded, under-resourced segregated schools for Black students. After the Board of Education refused to open the new John Philip Sousa Junior High School as an integrated facility, the group brought eleven Black students, including Spottswood Thomas Bolling Jr., to Sousa on opening day, September 11, 1950. Although the school had empty classrooms, the principal turned them away solely because of their race.

Howard law professors James M. Nabrit Jr. and George E. C. Hayes sued C. Melvin Sharpe, President of the Board, framing the case as a direct challenge to segregation itself rather than to unequal facilities. Because the District is a federal enclave, the Fourteenth Amendment's Equal Protection Clause, which binds only ``any State,'' did not apply; petitioners relied primarily on the Fifth Amendment's Due Process Clause.

\section{Procedural History}

The district court dismissed the complaint, relying on \textit{Carr v. Corning}, 182 F.2d 14 (D.C. Cir. 1950), which had upheld the District's dual school system. While the appeal was pending, the Supreme Court granted certiorari before judgment so the case could be heard alongside the four state cases in \textit{Brown v. Board of Education}. After argument in 1952 and reargument in 1953, the Court decided the case on May 17, 1954, the same day as \textit{Brown}, and restored it to the docket for argument on the remedy, later resolved in \textit{Brown II}.

\section{Judicial Votes}

Unanimous, \textbf{9--0}.
\begin{itemize}
\item \textbf{Majority:} Chief Justice Warren, joined by Justices Black, Reed, Frankfurter, Douglas, Jackson, Burton, Clark, and Minton.
\item \textbf{Concurrences and dissents:} None.
\end{itemize}

\section{Holding}

Racial segregation in the public schools of the District of Columbia violates the Due Process Clause of the Fifth Amendment. Although the Fifth Amendment lacks an explicit Equal Protection Clause, due process and equal protection are not mutually exclusive, and federal discrimination may be so unjustifiable as to violate due process. Racial classifications are constitutionally suspect; because school segregation is not reasonably related to any proper governmental objective, it arbitrarily deprives Black students of liberty. The strict holding is narrow; Warren's broader statements were expansive rationale that later decisions, culminating in \textit{Adarand Constructors, Inc. v. Pena} (1995), elevated into ``reverse incorporation.''

\section{Analysis of Opinions}

\textbf{The textual gap.} Warren's brief opinion conceded that the Fifth Amendment contains no equal protection clause and that the two guarantees are not ``always interchangeable phrases,'' but held them ``not mutually exclusive'' because both stem from ``our American ideal of fairness.'' Citing \textit{Detroit Bank}, \textit{Currin v. Wallace}, and \textit{Steward Machine Co.}, he built the bridge: discrimination may be so unjustifiable as to violate due process.

\textbf{Suspect classifications.} Relying on the federal wartime cases \textit{Korematsu} and \textit{Hi\-ra\-ba\-ya\-shi}, Warren reiterated that racial classifications are ``constitutionally suspect,'' and drew support from \textit{Gibson v. Mississippi} and \textit{Buchanan v. Warley}.

\textbf{Liberty.} Rejecting liberty as ``mere freedom from bodily restraint,'' Warren held that segregation, ``not reasonably related to any proper governmental objective,'' arbitrarily deprives Black children of liberty. He did not repeat \textit{Brown}'s social-science findings or apply a modern ``compelling interest'' test.

\textbf{Structural appeal.} Given \textit{Brown}, ``it would be unthinkable that the same Constitution would impose a lesser duty on the Federal Government.'' Unanimity let the Court speak with one voice nationwide.

\section{Examples: Future Applications}

\begin{enumerate}
\item \textbf{Same-Side (Invalidated): Sex-Based Survivor Benefits.} A federal Social Security provision pays survivor benefits to widows with minor children but not to widowers. The federal government cannot enforce a sex classification that would violate the Fourteenth Amendment if enacted by a state (as in \textit{Weinberger v. Wiesenfeld}).
\item \textbf{Same-Side (Invalidated): Federal Contracting Preferences.} A federal highway program rewards hiring minority-owned subcontractors without a compelling remedial interest or narrow tailoring. Under \textit{Adarand}, federal racial classifications face the same strict scrutiny as state ones.
\item \textbf{Opposite-Side (Upheld): Alienage Limits in Medicare.} Congress conditions a federal medical insurance program on five years' lawful residence for noncitizens. Given Congress's immigration power, \textit{Bolling}'s caveat that the guarantees are not ``always interchangeable'' permits rational-basis review (as in \textit{Mathews v. Diaz}).
\item \textbf{Opposite-Side (Upheld): Brewery Excise Taxes.} Congress taxes large breweries more heavily than craft breweries. With no suspect class or fundamental right, the classification easily satisfies Fifth Amendment due process.
\item \textbf{Fence-Sitter: Benefits in the Territories.} Congress excludes Puerto Rico residents from SSI. Challengers allege arbitrary geographic and ethnic discrimination; the government invokes the Territory Clause and Puerto Rico's tax status, a tension later resolved for deference in \textit{United States v. Vaello Madero} (2022).
\end{enumerate}

\section{Critique}

\textbf{Originalist and textualist views.} No one in 1791, when the Constitution accommodated slavery, understood the Fifth Amendment to bar federal racial discrimination, and reading equal protection into due process makes the Fourteenth Amendment's separate clause redundant. Michael McConnell argues the Court could have held the segregation policy ultra vires; Ryan Williams and Akhil Amar would ground the result in the Citizenship Clause.

\textbf{Progressive and structural views.} Critics call the ``unthinkable'' rationale ipse dixit, justified, in Richard Primus's words, ``by the force of sheer normative necessity.'' Its revival of broad substantive due process sat uneasily with the repudiation of \textit{Lochner}, and \textit{Adarand}'s rigid congruence ignored that the Reconstruction Framers distrusted state majorities more than Congress and empowered Congress to enact remedial race-conscious measures.

\section{Key Quotations}

\begin{quote}
``The `equal protection of the laws' is a more explicit safeguard of prohibited unfairness than `due process of law,' and therefore we do not imply that the two are always interchangeable phrases. But, as this Court has recognized, discrimination may be so unjustifiable as to be violative of due process.''
\end{quote}
\noindent --- Chief Justice Warren, Opinion of the Court, 347 U.S. at 499.

\begin{quote}
``The Fifth Amendment, which is applicable in the District of Columbia, does not contain an equal protection clause, as does the Fourteenth Amendment, which applies only to the states. But the concepts of equal protection and due process, both stemming from our American ideal of fairness, are not mutually exclusive.''
\end{quote}
\noindent --- Chief Justice Warren, Opinion of the Court, 347 U.S. at 499.

\begin{quote}
``Liberty under law extends to the full range of conduct which the individual is free to pursue, and it cannot be restricted except for a proper governmental objective. Segregation in public education is not reasonably related to any proper governmental objective, and thus it imposes on Negro children of the District of Columbia a burden that constitutes an arbitrary deprivation of their liberty in violation of the Due Process Clause.''
\end{quote}
\noindent --- Chief Justice Warren, Opinion of the Court, 347 U.S. at 499--500.

\end{document}
